\chapter{线性变换}
\orig{42--49}

\section{线性变换}

\subsection{定义与非线性变换}

\textbf{1. 定义。} 线性空间 $V$ 的一个变换 $A$ 称为\textbf{线性变换}，若对任意 $\alpha,\beta\in V$ 和任意 $k\in F$，都有
\[
A(\alpha+\beta)=A(\alpha)+A(\beta),
\qquad
A(k\alpha)=kA(\alpha).
\]

\textbf{2. 非线性变换的例。}

\textbf{例 1。} 取 $F=\mathbb R$，
\[
V=\{(x_1,x_2,x_3)\mid x_i\in F\},
\]
令
\[
A(x_1,x_2,x_3)=(x_1^2,x_2+x_3,x_3^2).
\]
当
\[
\alpha=(1,2,3),\qquad \beta=(1,1,1)
\]
时，
\[
A(\alpha)=(1,5,9),\qquad A(\beta)=(1,2,1),
\]
而
\[
A(\alpha+\beta)=A(2,3,4)=(4,7,16)\ne A(\alpha)+A(\beta),
\]
故不满足加法条件。另一方面
\[
A(k\alpha)=(k^2,5k,9k^2),
\qquad
kA(\alpha)=(k,5k,9k),
\]
当 $k\ne0,1$ 时一般不等，因此齐次条件也不成立。\jiaokan{所选数值例以及“取 $k\ne0,1$”的论证并不适用于只有两个元素的域。为避免无关的特征例外，此处明确取 $F=\mathbb R$；反例目的不变。}

\textbf{例 2。} 取 $F=\mathbb R$，$V=\mathbb R^2$，定义
\[
A(a,b)=\begin{cases}
(a,b),&a,b\text{ 同号，或至少有一个为 }0,\\
(-a,-b),&a,b\text{ 异号}.
\end{cases}
\]
容易验证 $A(k(a,b))=kA(a,b)$，即齐次条件成立；但
\[
A(-2,-3)=(-2,-3),\qquad A(-1,4)=(1,-4),
\]
而
\[
A((-2,-3)+(-1,4))=A(-3,1)=(3,-1)
\]
不等于
\[
(-2,-3)+(1,-4),
\]
所以加法条件不成立。这说明两个线性条件彼此独立。

\textbf{例 3。} 把复数域 $\mathbb C$ 看作它自身上的线性空间，令
\[
A(\alpha)=\bar\alpha.
\]
则
\[
A(\alpha+\beta)=\bar\alpha+\bar\beta=A(\alpha)+A(\beta),
\]
但取 $\alpha=1,k=i$，有
\[
A(k\alpha)=-i\ne i=kA(\alpha),
\]
所以它不是复线性变换。这再次说明齐次条件独立于加法条件。

判断一个变换不是线性变换，只需找到一条线性条件失败即可；上面三例同时把各种可能情形都展示出来。

\subsection{线性变换与一一变换}

\textbf{例 1。} 零变换是线性变换，但在非零线性空间中不是一一变换。

\textbf{例 2。} 复共轭 $A(\alpha)=\bar\alpha$ 是一一变换，但在 $\mathbb C$ 作为复线性空间时不是线性变换。

\textbf{例 3。} 恒等变换 $E(\alpha)=\alpha$ 既是线性变换，又是一一变换。

\subsection{若干反例}

\textbf{4.} 已知线性变换把线性相关的向量组变为线性相关的向量组，但反之不真。零变换就把线性无关的向量组变成全零向量组，因而变成线性相关组。

\textbf{5.} 在特征 $0$ 的数域上，不存在有限阶矩阵 $A,B$ 满足
\[
AB-BA=E,
\]
但在线性变换中可以存在这样的 $A,B$。在 $F[x]$ 上令
\[
Af(x)=f'(x),
\qquad
Bf(x)=xf(x).
\]
则
\[
(AB-BA)f(x)=A(xf(x))-x f'(x)=f(x),
\]
从而
\[
AB-BA=E.
\]
\jiaokan{有限矩阵“不可能满足 $AB-BA=E$”要借助迹，需有 $\operatorname{tr}E=n\ne0$；因此在一般正特征域中不能无条件表述。这里按“数域”通常指特征 $0$ 的语境补明。}

\textbf{6.} 线性变换的乘法不满足交换律。

例如在实数域上的多项式空间 $\mathbb R[x]$ 中定义
\[
D(f(x))=f'(x),
\qquad
J(f(x))=\int_0^x f(t)\,dt.
\]
则
\[
DJ=E,
\]
而一般
\[
JD\ne E.
\]

\textbf{7.} 相似矩阵有相同的特征多项式，但反之不真。

例如
\[
A=\begin{pmatrix}1&0\\0&1\end{pmatrix},
\qquad
B=\begin{pmatrix}1&1\\0&1\end{pmatrix}.
\]
二者都有
\[
|\lambda E-A|=|\lambda E-B|=(\lambda-1)^2,
\]
但 $A$ 与 $B$ 不相似，因为任何与 $A=E$ 相似的矩阵仍是 $E$，而 $B\ne E$。

\subsection{Hamilton--Cayley 定理的两个注意点}

设 $A$ 是数域 $F$ 上的一个 $n\times n$ 矩阵，
\[
f(\lambda)=|\lambda E-A|
\]
是 $A$ 的特征多项式，则
\[
f(A)=0.
\]
对此应注意两点。

1) $A$ 与特征值 $\lambda_0$ 都可以说是 $f$ 的“根”，但意义不同。特征值是某个扩域中的标量，而 $A$ 是矩阵环 $M_n(F)$ 中的元素；当 $n\ge2$ 时，矩阵环不是域，也不能作为 $F$ 的扩域来理解。

2) 在域中，一个 $n$ 次非零多项式至多有 $n$ 个根；在含零因子的环中则未必如此。

例如取 $F=\mathbb Q$，对任意 $a\in\mathbb Q$，令
\[
A_a=\begin{pmatrix}1&a\\0&1\end{pmatrix}.
\]
每个 $A_a$ 都满足
\[
f(\lambda)=|\lambda E-A_a|=(\lambda-1)^2=\lambda^2-2\lambda+1,
\]
并且
\[
f(A_a)=A_a^2-2A_a+E=0.
\]
因此一个二次多项式在矩阵环中可以有无穷多个“根”。

\section{线性变换的值域与核}

\textbf{1. 定义。} 若 $A$ 是线性空间 $V$ 的线性变换，$A$ 的所有像所成的集合称为 $A$ 的\textbf{值域}，记为 $AV$；在 $A$ 下零向量的所有逆像所成的集合称为 $A$ 的\textbf{核}，记为 $A^{-1}(0)$。

\textbf{2.} 已知
\[
\dim AV+\dim A^{-1}(0)=\dim V,
\]
但 $AV$ 与 $A^{-1}(0)$ 的和未必等于 $V$。

例如令 $F[x]_n$ 表示数域 $F$ 上所有次数不超过 $n-1$ 的多项式及零多项式所成的空间，定义
\[
Df(x)=f'(x).
\]
在特征 $0$ 的数域语境下，
\[
DV=F[x]_{n-1},
\qquad
D^{-1}(0)=F,
\]
且
\[
\dim F[x]_n=n,
\quad
\dim F[x]_{n-1}=n-1,
\quad
\dim F=1.
\]
但是
\[
F[x]_{n-1}+F=F[x]_{n-1}\ne F[x]_n,
\]
原因是
\[
F[x]_{n-1}\cap F=F\ne0.
\]

\section{不变子空间}

\textbf{1. 定义。} 设 $A$ 是数域 $F$ 上线性空间 $V$ 的线性变换，$S$ 是 $V$ 的子空间。若对任意 $\alpha\in S$ 都有 $A\alpha\in S$，则称 $S$ 是 $A$ 的\textbf{不变子空间}，简称 $A$--子空间。

\textbf{2.} 不变子空间是相对于线性变换而言的。同一子空间对一个变换可以不变，对另一个变换则未必不变。

例如取
\[
V=F^3,
\]
\[
A(x_1,x_2,x_3)=(2x_1-x_2,\ x_2+x_3,\ x_1),
\]
\[
B(x_1,x_2,x_3)=(0,x_2,x_3).
\]
显然 $A,B$ 都是线性变换，且
\[
BV=\{(0,x_2,x_3)\},
\qquad
B^{-1}(0)=\{(x_1,0,0)\}.
\]
这两个子空间都是 $B$--子空间，但都不是 $A$--子空间。例如
\[
(0,1,1)\in BV,
\qquad
A(0,1,1)=(-1,2,0)\notin BV,
\]
又
\[
(1,0,0)\in B^{-1}(0),
\qquad
A(1,0,0)=(2,0,1)\notin B^{-1}(0).
\]

此外
\[
AB(x_1,x_2,x_3)=(-x_2,x_2+x_3,0),
\]
而
\[
BA(x_1,x_2,x_3)=(0,x_2+x_3,x_1),
\]
所以 $AB\ne BA$。这个例子同时再次说明线性变换乘法不满足交换律。

值得注意的是，在这个例子中 $A$ 实际上是可逆变换。给定 $(y_1,y_2,y_3)\in V$，由
\[
2x_1-x_2=y_1,
\qquad
x_2+x_3=y_2,
\qquad
x_1=y_3
\]
可得
\[
(x_1,x_2,x_3)=(y_3,-y_1+2y_3,y_2+y_1-2y_3),
\]
故
\[
AV=V,
\qquad
A^{-1}(0)=0.
\]
于是 $AV$ 与 $A^{-1}(0)$ 当然都是 $B$--子空间。

\textbf{3.} 若 $A,B$ 是 $V$ 的线性变换并且 $AB=BA$，则 $BV$ 与 $B^{-1}(0)$ 都是 $A$--子空间；同样，$AV$ 与 $A^{-1}(0)$ 都是 $B$--子空间。但反之不真。

例如取
\[
A(x_1,x_2,x_3)=(x_1-x_2,x_2,x_3),
\qquad
B(x_1,x_2,x_3)=(-x_1,x_2,x_3).
\]
二者都是可逆线性变换，因而
\[
AV=BV=V,
\qquad
A^{-1}(0)=B^{-1}(0)=0;
\]
这些子空间对任意线性变换都不变。但是
\[
AB(x_1,x_2,x_3)=(-x_1-x_2,x_2,x_3),
\]
而
\[
BA(x_1,x_2,x_3)=(-x_1+x_2,x_2,x_3),
\]
故
\[
AB\ne BA.
\]
